\uuid{HDym}
\exo7id{1196}
\titre{exo7 1196}
\auteur{legall}
\organisation{exo7}
\datecreate{1998-09-01}
\isIndication{false}
\isCorrection{false}
\chapitre{Suite}
\sousChapitre{Convergence}
\cpge{mpsi-mp2i.nombres-reels-suites.c, pcsi.nombres-reels-suites.b, ptsi.nombres-reels-suites.b, mpsi-mp2i.nombres-reels-suites.d, pcsi.nombres-reels-suites.c, ptsi.nombres-reels-suites.c, mpsi-mp2i.nombres-reels-suites.e, pcsi.nombres-reels-suites.d, ptsi.nombres-reels-suites.d, bcpst1.suites-reelles}
\module{Analyse}
\niveau{L1}
\difficulte{}

\contenu{
\question{
Soient $  (u_n)_{n\in { \Nn}}  $ et $  (v_n)_{n\in { \Nn}}  $
deux suites de nombres r\' eels telles que $  0<u_1 <v_1  $ et $  u_{n+1}= \sqrt{ u_nv_n}  $
et $  \displaystyle{ v_{n+1}=\frac{u_n+v_n} { 2}}  .$ Montrer qu'elles convergent vers la m\^eme limite.
}
}
